% P15.21: Feller VII.7 Theorem 2 and Durrett Examples 2.2.15--2.2.16.
\chapter{Weak Normalization with Infinite Means}
\label{sec:p15-infinite-mean-weak-normalization}

A positive gain can have infinite expectation while the total from many
independent trials concentrates around a deterministic scale. The scale
then need not be proportional to the number of trials. This chapter
gives Feller's necessary and sufficient condition, constructs its
normalizers, and retains the separate triangular-array calculation for
the St.~Petersburg game. The sources are Feller, Volume II, second
edition, VII.7, Theorem~2 and its continuation
\cite[pp.~236--237]{Feller1971ProbabilityII}, and Durrett's author
Version~5 dated 11 January 2019, Examples~2.2.15--2.2.16
\cite[pp.~65--66]{DurrettPTEAuthor2019}.

\section{Positive gains and exact truncation conventions}
\label{sec:p15-positive-gain-conventions}

\begin{hbdefinition}{def:p15-positive-gain-normalization}{The law, its tail and its truncated moments}
Let \(X_1,X_2,\ldots\) be independent identically distributed finite
real random variables on \((\Omega,\mathcal F,\mathbb P)\), with
\(X_1\geq0\) almost surely and \(\mathbb P(X_1>0)>0\).
Write \(X\) for a variable with this common law,
\(F(s)=\mathbb P(X\leq s)\), and \(S_n=\sum_{k=1}^nX_k\).
Feller's stated domain is the special case \(\mathbb P(X>0)=1\),
equivalently \(F(0)=0\). The results below also prove the indicated
extension allowing an atom at zero.

For \(s>0\), define
\[
 q(s)=\mathbb P(X>s),\quad
 \mu(s)=\mathbb E[X\mathbf1_{\{X\leq s\}}],\quad
 v(s)=\mathbb E[X^2\mathbf1_{\{X\leq s\}}],
\]
and
\[
 \mu_<(s)=\mathbb E[X\mathbf1_{\{X<s\}}],\qquad
 q_{\geq}(s)=\mathbb P(X\geq s).
\]
These moments are finite even when \(\mathbb EX=\infty\).
Choose any \(s_0>0\) with \(\mu(s_0)>0\). For \(s\geq s_0\), put
\[
 \delta(s)=\frac{s q(s)}{\mu(s)},\qquad
 \rho(s)=\frac{\mu(s)}{s q(s)}.
\]
If \(q(s)=0\), set \(\rho(s)=+\infty\) and \(\delta(s)=0\).
Thus \(\rho(s)\to\infty\) means exactly \(\delta(s)\to0\).
No ratio with zero truncated mean is used.
A \emph{relative weak normalizer} is a deterministic sequence
\(a_n>0\) such that \(S_n/a_n\to1\) in probability:
\(\mathbb P(|S_n/a_n-1|>\varepsilon)\to0\) for every
\(\varepsilon>0\).
\index{weak law!infinite mean}\index{normalization!relative weak}
\end{hbdefinition}

\begin{hblemma}{thm:p15-positive-truncation-tools}{proved=annex:p15.positive-truncation-tools}{Truncated moment identities and the two-error bound}
For every \(s>0\),
\[
 \mu(s)=\int_0^s q(t)\,dt-sq(s),\qquad
 v(s)=2\int_0^s t q(t)\,dt-s^2q(s),\qquad
 v(s)\leq s\mu(s).
\]
For \(a>0\), \(\varepsilon>0\) and \(n\geq1\),
\[
 \mathbb P(|S_n-n\mu(s)|>\varepsilon a)
 \leq \frac{n v(s)}{\varepsilon^2a^2}+n q(s).
\]
Moreover \(\mu(s)/s\to0\). If \(\delta(s)\to0\), then
\[
 \frac{v(s)}{s\mu(s)}\longrightarrow0.
\]
The deletion truncation retains the atom at \(s\); the error event
uses the strict tail \(X>s\).
\end{hblemma}

\begin{hblemma}{thm:p15-relative-normalizer-properties}{proved=annex:p15.relative-normalizer-properties}{Growth, adjacent ratios and uniqueness of scale}
Under the positive-gain assumptions, every relative weak normalizer
satisfies \(a_n\geq c n\) for all sufficiently large \(n\), for some
\(c>0\), and
\[
 a_n\longrightarrow\infty,\qquad
 \frac{a_{n+1}}{a_n}\longrightarrow1.
\]
No monotonicity of \(a_n\) is assumed. If \(a_n\) and \(b_n\) both
normalize \(S_n\) relatively to one, then \(a_n/b_n\to1\).
Conversely, replacing a relative weak normalizer by positive
\(b_n\) with \(b_n/a_n\to1\) preserves the relative weak limit.
\end{hblemma}

\section{The normalizer construction and Feller's criterion}
\label{sec:p15-feller-positive-criterion}

\begin{hblemma}{thm:p15-positive-normalizer-construction}{proved=annex:p15.positive-normalizer-construction}{An exact anchored crossing}
Choose \(s_0>0\) with \(\mu(s_0)>0\), and choose an integer \(N_0\)
so that \(n\mu(s_0)>s_0\) for \(n\geq N_0\).
For those integers define
\[
 s_n=\inf\{s\geq s_0:n\mu(s)/s\leq1\}.
\]
The set is nonempty, its infimum is finite and exceeds \(s_0\), and
\[
 n\mu(s_n)=s_n,\qquad
 s_n\geq n\mu(s_0)\longrightarrow\infty.
\]
The sequence \(s_n\), \(n\geq N_0\), is nondecreasing.
These existence statements require no finite expectation and no
assumption on \(\delta\). Arbitrary positive values may be assigned
at the finitely many indices below \(N_0\).
\end{hblemma}

\begin{hbdefinition}{def:p15-positive-slow-variation}{Slow variation of the truncated mean}
An eventually positive function \(L:(0,\infty)\to[0,\infty)\)
is \emph{slowly varying at infinity} if
\(L(cs)/L(s)\to1\) as \(s\to\infty\) for every fixed \(c>0\).
Here this definition is applied to the nondecreasing truncated
mean \(\mu\), with its eventual positive domain understood.
\index{slow variation!truncated mean}
\end{hbdefinition}

\begin{hbtheorem}{thm:p15-positive-tail-equivalences}{proved=annex:p15.positive-tail-equivalences}{Slow variation and either endpoint convention}
Under the positive-gain assumptions, the following are equivalent:
\[
 \delta(s)\to0;\qquad
 \mu\text{ is slowly varying};\qquad
 \frac{s q_{\geq}(s)}{\mu_<(s)}\to0.
\]
The last ratio is taken only on the eventual domain
\(\mu_<(s)>0\). Under these conditions,
\[
 \frac{\mu_<(s)}{\mu(s)}\to1,\qquad
 \frac{s\,\mathbb P(X=s)}{\mu(s)}\to0.
\]
In the infinite-mean case this is exactly the truncated-first-moment
specialization of the relation noted in Feller's footnote to VII.7,
Theorem~2; the needed equivalence is proved here directly
\cite[pp.~236, 282--284]{Feller1971ProbabilityII}.
The finite-mean case is included as well.
\end{hbtheorem}

\begin{hbtheorem}{thm:p15-feller-positive-weak-normalization}{proved=annex:p15.feller-positive-weak-normalization}{Feller VII.7, Theorem 2: both directions}
For independent identically distributed nonnegative gains with
\(\mathbb P(X>0)>0\), there exist positive deterministic \(a_n\)
with \(S_n/a_n\to1\) in probability if and only if
\[
 \rho(s)\longrightarrow\infty\quad(s\to\infty).
\]
Allowing arbitrary nonzero real constants does not enlarge the
conclusion: any sequence giving this limit is eventually positive.
When the condition holds, the anchored construction gives
\[
 a_n=n\mu(s_n)=s_n\qquad(n\geq N_0)
\]
and, for every \(\varepsilon>0\),
\[
 \mathbb P(|S_n/s_n-1|>\varepsilon)
 \leq \frac{v(s_n)}{\varepsilon^2s_n\mu(s_n)}
                              +\delta(s_n)\longrightarrow0.
\]
Feller's sufficiency argument uses the truncated second moment.
His necessity argument uses positivity, the largest positive
summand, and truncation at \(2a_n\). Both arguments, including the
passage from the sequence \(2a_n\) to every large real threshold,
are supplied in the annex.
\end{hbtheorem}

\begin{hbcorollary}{thm:p15-positive-normalized-centers}{proved=annex:p15.positive-normalized-centers}{All relative normalizers and their truncated centers}
Suppose the equivalent conditions of the preceding theorem hold,
and let \(s_n\) be its constructed normalizer. A positive sequence
\(a_n\) is a relative weak normalizer if and only if \(a_n/s_n\to1\).
For every such \(a_n\),
\[
 \frac{n\mu(a_n)}{a_n}\to1,\quad
 \frac{n\mu_<(a_n)}{a_n}\to1,\quad
 n q(a_n)\to0,\quad n q_{\geq}(a_n)\to0.
\]
Thus the closed and strict truncated centers have the same leading
scale, although only the closed mean is used in the exact
right-continuous crossing construction.
If \(0<\mathbb EX=m<\infty\), the canonical normalizer satisfies
\(s_n/n\to m\), and \(nm\) is a relative weak normalizer.
If \(\mathbb EX=\infty\), then \(s_n/n\to\infty\).
\end{hbcorollary}

\section{Worked positive laws and required hypotheses}
\label{sec:p15-positive-normalization-examples}

\begin{hbtheorem}{thm:p15-positive-normalization-boundaries}{proved=annex:p15.positive-normalization-boundaries}{Finite means, a critical tail and failures}
\begin{enumerate}
\item If \(X=c>0\) almost surely and \(s_0=c\), then \(s_n=nc\)
for \(n\geq2\). If \(X\) equals \(1\) or \(3\), each with
probability \(1/2\), the choice \(s_0=3\) gives \(s_n=2n\)
for \(n\geq2\).
\item For the probability density
\(f(x)=x^{-2}\mathbf1_{\{x\geq1\}}\),
\[
 q(s)=s^{-1},\quad \mu(s)=\log s,\quad
 v(s)=s-1\qquad(s\geq1).
\]
Its mean is infinite and its relative weak normalizers satisfy
\(s_n\sim n\log n\), where these logarithms are natural.
\item For \(0<\alpha<1\), the probability density
\(f_\alpha(x)=\alpha x^{-\alpha-1}\mathbf1_{\{x\geq1\}}\)
has infinite mean and
\[
 \mu(s)=\frac{\alpha}{1-\alpha}(s^{1-\alpha}-1),\qquad
 \rho(s)\longrightarrow\frac{\alpha}{1-\alpha}.
\]
It therefore has no relative weak normalizer.
\item If all \(X_k\) are the same variable taking values \(1,3\)
with equal probabilities, their marginal law meets the tail
criterion, but no relative weak normalizer exists.
The identically zero law likewise has no normalizer giving limit one;
it is excluded by \(\mathbb P(X>0)>0\).
\end{enumerate}
Here \(u_n\sim v_n\) means \(u_n/v_n\to1\) for positive quantities.
\end{hbtheorem}

\section{The Cauchy boundary and St. Petersburg}
\label{sec:p15-cauchy-and-petersburg}

Each of the following source examples specifies its own common law
and reuses \(S_n=\sum_{k=1}^nX_k\) for that law. The Cauchy variables
are signed; the positive-gain hypotheses apply again to
St.~Petersburg variables.

\begin{hbdefinition}{def:p15-cauchy-petersburg-laws}{The two source laws}
For \(a>0\), the centered Cauchy law with scale \(a\) has density
\[
 f_a(x)=\frac{a}{\pi(a^2+x^2)}\qquad(x\in\mathbb R).
\]
The standard Cauchy law has \(a=1\).
The St.~Petersburg gain law has masses
\[
 \mathbb P(X=2^j)=2^{-j},\qquad j=1,2,\ldots.
\]
It is realized by independent fair coin tosses, paying \(2^j\) if the
first head is on toss \(j\). The event of no head has probability zero.
Repeated plays use independent copies of the gain.
In the St.~Petersburg statements below, \(\log_2\) always denotes
base-two logarithm.
\index{Cauchy distribution}\index{St. Petersburg paradox}
\end{hbdefinition}

\begin{hbtheorem}{thm:p15-cauchy-weak-law-failure}{proved=annex:p15.cauchy-weak-law-failure}{Durrett Example 2.2.15 and its distributional assertion}
The Cauchy densities integrate to one. For a standard Cauchy \(X\),
\[
 \mathbb P(|X|>x)\sim\frac{2}{\pi x}\quad(x\to\infty),
 \qquad \mathbb E|X|=\infty.
\]
For independent standard Cauchy \(X_k\), no real deterministic
sequence \(c_n\) satisfies \(S_n/n-c_n\to0\) in probability.
In fact \(S_n/n\) has exactly the standard Cauchy law for every
\(n\geq1\). More generally \(f_a*f_b=f_{a+b}\) almost everywhere
as densities for \(a,b>0\).
The tail argument is Durrett's proof in Example~2.2.15.
The annex additionally proves its later distributional assertion by
direct real convolution, without a characteristic-function premise.
\end{hbtheorem}

\begin{hbtheorem}{thm:p15-petersburg-array-calculation}{proved=annex:p15.petersburg-array-calculation}{Durrett Example 2.2.16: the complete truncation route}
Let \(X_k\) be independent St.~Petersburg gains.
Their common expectation is infinite. Choose real \(K(n)\to\infty\)
so that \(m(n)=\log_2 n+K(n)\) is an integer, and put
\[
 b_n=2^{m(n)}=n2^{K(n)},\qquad
 \overline X_{n,k}=X_k\mathbf1_{\{X_k\leq b_n\}},\qquad
 a_n=nm(n).
\]
For all sufficiently large \(n\), \(m(n)\geq1\), and
\[
 \begin{aligned}
 n\mathbb P(X>b_n)&=2^{-K(n)},&
 n\mathbb P(X\geq b_n)&=2^{1-K(n)},\\
 \mathbb E\overline X_{n,k}&=m(n),&
 \mathbb E\overline X_{n,k}^{\,2}&=2^{m(n)+1}-2\leq2b_n.
 \end{aligned}
\]
Consequently \((S_n-a_n)/b_n\to0\) in probability.
Whenever \(K(n)/\log_2n\to0\), the deterministic center satisfies
\(a_n/(n\log_2n)\to1\).
If, in addition, \(K(n)\leq\log_2\log_2n\) eventually, then
\[
 \frac{S_n}{n\log_2n}\longrightarrow1
                                      \quad\hbox{in probability}.
\]
All these conditions are simultaneously met, for \(n\geq2\), by
\[
 m(n)=\left\lfloor\log_2(n\log_2n)\right\rfloor,\qquad
 K(n)=m(n)-\log_2n.
\]
The floor \(\lfloor u\rfloor\) is the greatest integer at most \(u\).
Initial finitely many choices do not affect any assertion.

The cutoff follows Durrett's principle for applying
\cref{thm:p15-triangular-array-weak-law}: in checking its second
condition, \(w_n\to0\), the level \(b_n\) should be as small as possible
while the first condition, \(u_n\to0\), still holds
\cite[p.~65]{DurrettPTEAuthor2019}; the conclusion is obtained at the
scale \(b_n\). For this law the first condition alone fixes the order
of the smallest admissible levels: \(n\mathbb P(X>b)\geq n/b\) for every
\(b\geq1\), so \(u_n=n\mathbb P(X>b_n)\to0\) forces \(b_n/n\to\infty\).
For the dyadic levels \(b_n=n2^{K(n)}\) both conditions hold exactly
when \(K(n)\to\infty\), since \(u_n=2^{-K(n)}\) and
\(w_n\leq2^{1-K(n)}\); the second condition imposes nothing further.
\end{hbtheorem}

\begin{hbtheorem}{thm:p15-petersburg-normalizers-and-prices}{proved=annex:p15.petersburg-normalizers-and-prices}{Feller's route, finite calculations and the meaning of the scale}
For St.~Petersburg gains and \(s\geq2\), put
\(\ell(s)=\lfloor\log_2s\rfloor\). Then
\[
 \mu(s)=\ell(s),\quad
 v(s)=2^{\ell(s)+1}-2,\quad
 q(s)=2^{-\ell(s)},\quad
 \delta(s)\leq\frac{2}{\ell(s)}\longrightarrow0.
\]
Thus Feller's separate construction with \(s_0=2\) gives exact
normalizers \(s_n=n\lfloor\log_2s_n\rfloor\), for all sufficiently
large \(n\), and \(s_n/(n\log_2n)\to1\).
For the explicit array choice in the preceding theorem, \(n=1024\)
gives \(m(n)=13\), \(b_n=8192\), \(a_n=13312\), and
\[
 n\mathbb P(X>b_n)=\tfrac18,\qquad
 \frac{n\mathbb E\overline X_{n,k}^{\,2}}{b_n^2}
                   =\frac14-\frac1{32768}.
\]
The relative weak limit assigns the leading aggregate scale
\(n\log_2n\), or \(\log_2n\) per play. At \(n=1024\) this
per-play scale is \(10\). The scale equals \(40\) when
\(n=2^{40}=1\,099\,511\,627\,776\).

Both sources read this scale as a price. Durrett states the paradox:
the expectation of the gain is infinite, yet nobody would pay an
infinite amount to play. He presents the triangular-array theorem as
saying how much one should pay to play \(n\) times, namely a fair price
of \(\log_2n\) dollars per play, which is \(10\) dollars at
\(n=1024\). He then recalls Nicolas Bernoulli's remark of 1713 that
anyone even moderately sensible would sell the right to play for
\(40\) ducats per play, and computes that with a one-ducat wager about
\(2^{40}\approx10^{12}\) plays would be needed to start to break even
\cite[pp.~65--66]{DurrettPTEAuthor2019}. Feller introduces his
Theorem~2 through the theory of fair games discussed in his Volume~I,
Chapter~X: even when expectations exist, a participant in a fair game
can be strongly on the losing side, while the St.~Petersburg game shows
that the classical theory extends to some games of infinite expectation
provided the fair entrance fee depends on the number of trials. The
theorem makes this precise, with \(X_k\) read as the gains and \(a_n\)
as the total entrance fee for \(n\) trials
\cite[p.~236]{Feller1971ProbabilityII}.
In the sense proved here, a fee \(a_n\) is fair when
\(S_n/a_n\to1\) in probability. The fee is a scale for the total gain,
not an expectation, which remains infinite, and the theorem specifies
no finite-sample probability of breaking even. The \(2^{40}\) count is
the value of \(n\) at which the per-play scale \(\log_2n\) reaches
\(40\). For every fixed fee \(c>0\) per play,
\(\mathbb P(S_n>cn)\to1\).

No deterministic \(c_n\) makes \(S_n/n-c_n\to0\) in probability:
at \(x=2^j\), \(x\mathbb P(X>x)=1\).
Changing the denominator is essential.
\end{hbtheorem}

\section{No nonzero almost-sure limit when the mean is infinite}
\label{sec:p15-w02c-p15-018-infinite-mean-limits}

The relative weak normalizers above describe convergence in
probability. The almost-sure behaviour at the same scales is governed
by Feller's tail-series dichotomy, \cref{thm:p15-feller-infinite-mean},
for independent identically distributed real \(X_1,X_2,\ldots\), with
\(S_n=\sum_{k=1}^nX_k\), \(\mathbb E|X_1|=\infty\), a Feller
normalization \(a_n\) and its threshold series \(Q\) from
\cref{def:p15-feller-normalizers}.

Durrett introduces his Section~2.5.2 with the St.~Petersburg game,
for which \(S_n/(n\log_2n)\to1\) in probability while the limit
superior of the same ratio is almost surely infinite. He attributes the
theorem stated here as \cref{thm:p15-feller-infinite-mean} to Feller
(1946), reads it as showing that \(S_n/a_n\) cannot converge almost
surely to a nonzero limit when \(\mathbb E|X_1|=\infty\), and observes
that his earlier Theorem~2.3.8 was the special case \(a_n=n\)
\cite[p.~88]{DurrettPTEAuthor2019}; the theorem itself is his
Theorem~2.5.13 \cite[p.~89]{DurrettPTEAuthor2019}. The three results
below prove these consequences. The St.~Petersburg statement is
Durrett's Exercise~2.3.20 \cite[p.~76]{DurrettPTEAuthor2019}, and
Theorem~2.3.8 is \cite[p.~70]{DurrettPTEAuthor2019}.

\begin{hbtheorem}{thm:p15-w02c-p15-018-no-nonzero-limit}{proved=annex:p15-w02c-p15-018-no-nonzero-limit}{No nonzero almost-sure limit}
Let \(X_1,X_2,\ldots\) be i.i.d.\ real random variables with
\(\mathbb E|X_1|=\infty\), let \(a_n\) be a Feller normalization and
let \(Q=\sum_{n\geq1}\mathbb P(|X_1|\geq a_n)\). Exactly one of the
following holds:
\begin{enumerate}
\item \(Q<\infty\) and \(S_n/a_n\to0\) almost surely;
\item \(Q=\infty\) and \(\limsup_n|S_n|/a_n=\infty\) almost surely.
\end{enumerate}
Consequently the set of outcomes at which \(S_n/a_n\) converges to a
finite real limit different from zero is an event of probability zero.
This includes limits that vary with the outcome. If
\(\mathbb P(S_n/a_n\to c)>0\) for a real constant \(c\), then \(c=0\)
and \(Q<\infty\).
\end{hbtheorem}

\begin{hbcorollary}{thm:p15-w02c-p15-018-linear-scale}{proved=annex:p15-w02c-p15-018-linear-scale}{The linear scale: Durrett's Theorem 2.3.8}
Let \(X_1,X_2,\ldots\) be i.i.d.\ real random variables with
\(\mathbb E|X_1|=\infty\). Then
\[
 \sum_{n\geq1}\mathbb P(|X_1|\geq n)=\infty,\qquad
 \mathbb P(|X_n|\geq n\text{ i.o.})=1,
\]
and
\[
 \limsup_{n\to\infty}\frac{|S_n|}{n}=\infty\quad\text{almost surely}.
\]
In particular \(\mathbb P(\lim_nS_n/n\text{ exists in }\mathbb R)=0\).
The infinitely-often assertion and this last consequence are the two
conclusions of Durrett's Theorem~2.3.8
\cite[p.~70]{DurrettPTEAuthor2019}. The constant ratio
\(a_n/n=1\) makes \(a_n=n\) a Feller normalization, so this is the
linear case of the dichotomy. \Cref{thm:p15-power-rate-necessity} with
\(p=1\) gives the same moment conclusion from the opposite direction:
a zero rate \(S_n^{(c)}/n\to0\) forces \(\mathbb E|X_1|<\infty\).
\end{hbcorollary}

\begin{hbtheorem}{thm:p15-w02c-p15-018-petersburg-almost-sure}{proved=annex:p15-w02c-p15-018-petersburg}{St.~Petersburg: convergence in probability without almost-sure convergence}
Let \(X_1,X_2,\ldots\) be independent St.~Petersburg gains, with
base-two logarithms. For every real \(x>0\),
\[
 \mathbb P(X_1\geq x)\geq\min\{1,1/x\}.
\]
Almost surely,
\[
 \limsup_{n\to\infty}\frac{X_n}{n\log_2n}=\infty,\qquad
 \limsup_{n\to\infty}\frac{S_n}{n\log_2n}=\infty.
\]
Hence \(S_n/(n\log_2n)\to1\) in probability, by
\cref{thm:p15-petersburg-array-calculation}, but
\(\mathbb P\bigl(S_n/(n\log_2n)\to1\bigr)=0\).
With \(a_1=1\) and \(a_n=n\log_2n\) for \(n\geq2\), the sequence
\(a_n\) is a Feller normalization whose threshold series diverges, so
the second limit superior is also the case \(Q=\infty\) of
\cref{thm:p15-w02c-p15-018-no-nonzero-limit}.
\end{hbtheorem}

\begin{hbexample}{def:p15-w02c-p15-018-dyadic-ranges}{The threshold series in dyadic ranges}
Take St.~Petersburg gains and the threshold \(x=n\log_2n\).
At \(n=1024\), \(x=10\,240\). Since
\(2^{13}=8192<10\,240\leq16\,384=2^{14}\), the event
\(X_n\geq10\,240\) is \(X_n\geq2^{14}\), and
\[
 \mathbb P(X_n\geq10\,240)=\sum_{j\geq14}2^{-j}=2^{-13}
   =\frac1{8192}>\frac1{10\,240}.
\]
For the range \(1024\leq n\leq2047\), every index has
\(n\log_2n<2048\cdot11\), so the range contributes more than
\(1024/(2048\cdot11)=1/22\) to
\(\sum_n\mathbb P(X_n\geq n\log_2n)\). In general the range
\(2^k\leq n<2^{k+1}\) contributes more than \(1/(2(k+1))\), and these
lower bounds form a divergent harmonic tail. Each exceedance probability
is small, but their sum diverges, and the second Borel--Cantelli lemma
turns that divergence into infinitely many exceedances.
\end{hbexample}

\begin{hbwarning}{def:p15-w02c-p15-018-scope}{What the dichotomy does and does not exclude}
\begin{enumerate}
\item The infinite absolute mean is needed. With \(X_n=1\) and
\(a_n=n\), \(S_n/a_n=1\) for every \(n\)
(\cref{thm:p15-rate-hypothesis-boundaries}, item~4).
\item Independence is needed. If every \(X_n\) equals one positive
variable \(X\) with \(\mathbb P(X=2^j)=2^{-j}\), then
\(\mathbb E|X_1|=\infty\) and \(S_n/n=X\) for every \(n\): the ratio
converges at every outcome to a finite nonzero random limit
(\cref{thm:p15-rate-hypothesis-boundaries}, item~5).
\item The ratio condition on \(a_n/n\) is needed for the case split by
\(Q\): \cref{thm:p15-feller-ratio-counterexample} gives nondecreasing
denominators with \(Q<\infty\) and an infinite limit superior. This
section proves the no-nonzero-limit conclusion only for Feller
normalizations.
\item The dichotomy concerns limit superiors. For St.~Petersburg gains it
gives no information about \(\liminf_nS_n/(n\log_2n)\), and none is
asserted here.
\end{enumerate}
\end{hbwarning}

\emph{Relationships.} \Cref{thm:p15-w02c-p15-018-petersburg-almost-sure}
sharpens item~4 of \cref{thm:p15-rate-calculated-examples}, which has
the same dyadic law and gives \(\limsup_nS_n/n=\infty\); the denominator
\(n\log_2n\) here is the relative weak normalizer of
\cref{thm:p15-petersburg-array-calculation,thm:p15-petersburg-normalizers-and-prices}.
The finite-variance rate and its iterated-logarithm comparison are
\cref{thm:p15-finite-variance-log-rate} and
\cref{sec:p15-w02c-p15-018-rate-comparison}, and the parameter ranges
below order two are \cref{thm:p15-mz-full-classification}.

\begin{proofboundary}{def:p15-positive-normalization-boundary}{Exact continuation}
This chapter completes Feller VII.7, Theorem~2, its required
slow-variation equivalence, and Durrett Examples~2.2.15--2.2.16.
The general regular-variation theorem of Feller VIII.9 has wider
moment orders and endpoint cases than the specialization proved here.
Its full presentation, the alternate argument in XVII.2a, and the
separately identified Poisson/Bessel example remain separate source
components. The complete numbered weak-convergence presentation is
\cref{sec:p15-weak-convergence-of-laws}, including its quantile,
portmanteau, selection and metric arguments.
The distinct characteristic-function proof of the Cauchy transform
and average law is supplied in
\cref{thm:p15-characteristic-inverse-examples}.
The direct convolution proof here retains its separate route.
The complete remaining characteristic-function source chain is
\cref{sec:p15-characteristic-function-criteria-and-moments}.
The central-limit constructions and exercises are supplied in
\cref{sec:p15-central-limit-theorems}. The complete normal-domain
criterion is \cref{thm:p15-nda-characterization}, with its original
truncated-mean cutoff presentation in \cref{thm:p15-nda-gk-cutoff-route}.
The later full stable-law presentations retain their distinct source
obligations.
\end{proofboundary}
